% ============================================================================
% Section IV: Proposed LLMBO-MO
% Revised to emphasize the core methodological idea and reduce implementation-level detail.
% ============================================================================

\section{Proposed LLMBO-MO}
\label{sec:method}

\subsection{Framework}
\label{sec:framework}

LLMBO-MO uses battery-domain knowledge to guide the choice of protocols to
simulate. Figure~\ref{fig:framework} shows the optimization loop. The LLM
first proposes an initial set of protocols. During the early BO iterations,
it also suggests a region for the current trade-off among objectives. The BO
optimizer scalarizes the observed objectives, fits a GP, and selects the next
protocol by maximizing EI. Evaluating that protocol with the battery model
provides its objective values and feasibility, which are used to update the
data set $D$ and the feasible Pareto archive $P$. The loop stops when the
simulation budget is exhausted.

\begin{figure}[!t]
\centering
\includegraphics[width=\columnwidth]{fig2.pdf}
\caption{Structure of the proposed LLMBO-MO framework}
\label{fig:framework}
\end{figure}


\subsection{Decomposition-Based Multiobjective BO}
\label{sec:parego_backbone}

The numerical optimizer follows the decomposition-based approach of
ParEGO~\cite{ref4}. A set $W$ of nonnegative weight vectors represents
trade-offs among the three objectives. At iteration $t$, the optimizer selects
$\mathbf w^{(t)}=[w_1^{(t)},w_2^{(t)},w_3^{(t)}]$ from $W$, with
$\sum_{i=1}^{3}w_i^{(t)}=1$.

Charging time and aging are first transformed to reduce differences in scale:
\begin{equation}
 y_1(\theta)=\log\!\frac{t_c(\theta)}{t_{\mathrm{ref}}},
 y_2(\theta)=T_p(\theta),
 y_3(\theta)=\log\!\frac{Q_s(\theta)}{Q_{\mathrm{ref}}}
 \label{eq:objective_transform}
\end{equation}
where $t_{\mathrm{ref}}$ and $Q_{\mathrm{ref}}$ are fixed positive reference
values. A positive floor is applied before taking logarithms. Let
$z_{i,t}^{*}$ be the best observed value of transformed objective $i$ and
$s_{i,t}$ its scaling factor. The normalized distance is
\begin{equation}
 \bar f_{i,t}(\theta)=
 \frac{|y_i(\theta)-z_{i,t}^{*}|}{s_{i,t}},\qquad i=1,2,3.
 \label{eq:objective_scaling}
\end{equation}
The augmented Tchebycheff scalarization combines these distances:
\begin{equation}
 g_t(\theta)=
 \max_{i=1,2,3}\!\left\{w_i^{(t)}\bar f_{i,t}(\theta)\right\}
 +\rho_{\mathrm{aug}}\sum_{i=1}^{3}w_i^{(t)}\bar f_{i,t}(\theta)
 \label{eq:tchebycheff}
\end{equation}
where $\rho_{\mathrm{aug}}>0$ is a small augmentation coefficient. Changing
$\mathbf w^{(t)}$ directs successive iterations toward different parts of the
Pareto front.

A GP with an automatic relevance determination (ARD) Mat\'ern-$5/2$ kernel is
fitted to the scalarized observations. Its posterior mean and standard
deviation are denoted by $\mu_t(\theta)$ and $\sigma_t(\theta)$. For
minimization, EI is
\begin{align}
 \mathrm{EI}_t(\theta)
 &=\bigl(g_t^{\mathrm{best}}-\mu_t(\theta)\bigr)\Phi(x_t)
 +\sigma_t(\theta)\phi(x_t)\\
 x_t&=\frac{g_t^{\mathrm{best}}-\mu_t(\theta)}{\sigma_t(\theta)}
 \label{eq:ei}
\end{align}
where $g_t^{\mathrm{best}}$ is the best observed scalarized value, and
$\Phi$ and $\phi$ are the standard normal cumulative distribution and density
functions.

\subsection{LLM-Informed Initialization}
\label{sec:warm_start}

The first GP is fitted to a small initial data set. To make these evaluations
more informative for charging design, the LLM is given the battery context,
MCC protocol structure, decision-variable bounds, objectives, and domain
knowledge about charging (Fig.~\ref{fig:prompt_warmstart}). It returns
numerical candidate protocols and a preference ranking.

% Prompt exhibits for the two LLM touchpoints.
% Typeset in the style of boxed prompt figures: gray panel, thin frame,
% domain-knowledge phrases highlighted in red.
\newcommand{\promptbox}[1]{%
  \fcolorbox{black}{gray!10}{\parbox{\dimexpr\linewidth-2\fboxsep-2\fboxrule\relax}{#1}}}

\begin{figure}[!t]
\centering
\promptbox{\footnotesize
You are an expert in \textcolor{red}{lithium-ion battery fast-charging
optimization}. Cell: LG INR21700-M50 (5~Ah, NMC811/graphite); the charging
window is 0\% to 80\% SOC.\\
Decision variables and bounds: $I_1\in[2.0,6.0]$~A,
$I_2\in[2.0,5.0]$~A, $I_3\in[2.0,3.0]$~A,
$\Delta s_1\in[0.10,0.40]$, $\Delta s_2\in[0.10,0.30]$.\\
Constraints: keep $\Delta s_1+\Delta s_2\le 0.65$ whenever possible
(practical hard limit $0.70$).
\textcolor{red}{A decreasing current profile $I_1\ge I_2\ge I_3$ is
recommended, since high currents at low SOC carry a lower risk of lithium
plating than the same currents near full charge.}\\
Objectives: minimize charging time, peak temperature rise, and capacity
fade.\\
Generate $15$ \textcolor{red}{diverse} protocols that
\textcolor{red}{cover multiple trade-off directions instead of clustering
around one style}. Respond with only a JSON array.}
\caption{Prompt for LLM-informed warm-start initialization.}
\label{fig:prompt_warmstart}
\end{figure}

\begin{figure}[!t]
\centering
\promptbox{\footnotesize
You are guiding a battery charging Bayesian optimization loop.\\
Iteration $t$ of $T$; weight vector
$\mathbf w^{(t)}=(w_1,w_2,w_3)$; current best scalarized value
$g_t^{\mathrm{best}}$. Decision variables and bounds as in
Fig.~\ref{fig:prompt_warmstart}; hard constraint
$\Delta s_1+\Delta s_2<0.70$.\\
Return exactly one JSON object with one of two structures:\\
\hspace*{1em}\texttt{"point"}: \texttt{\{I1, I2, I3, dSOC1, dSOC2,
confidence\}}\\
\hspace*{1em}\texttt{"region"}: \texttt{\{lb: \{...\}, ub: \{...\},
confidence\}}\\
Guidelines: \textcolor{red}{use a point when confident about a specific
location; use a region when the promising area is uncertain} (region width
$5$--$20\%$ of the range per dimension). The confidence lies in $[0,1]$ and
\textcolor{red}{reflects your certainty}. All values must respect the
parameter bounds and the $\Delta s$ constraint. Output JSON only.}
\caption{Prompt for iteration-level region-preference guidance.}
\label{fig:prompt_region}
\end{figure}



The candidate pool is screened before simulation. Entries with non-numerical
or non-finite values, duplicates, and invalid protocol representations are
removed; out-of-bound values are repaired. The remaining candidates are
selected using their ranking and diversity in the decision space. Random
valid protocols supplement this selection to cover other parts of the search
space. The resulting $n_0$ protocols are simulated to form the initial data
set and Pareto archive.

\subsection{Region-Lift Guidance}
\label{sec:region_pref}

Initialization introduces LLM suggestions before sequential optimization
begins. During the search, the relevant trade-off changes with the scalarization
weights. At an early iteration $t$, the LLM therefore receives the current
weights $\mathbf w^{(t)}$ and the observed optimization state, and returns a
preferred region $R_t$ for that subproblem
(Fig.~\ref{fig:prompt_region}). Restricting the search to $R_t$
would exclude all candidates outside the proposed region. Instead, Region-Lift
uses GP posterior cross-covariance to adjust their acquisition values. This
extends LGBO's region-lifted preference mechanism~\cite{ref5} to the current
scalarized objective.

A valid region is represented by $M$ Sobol anchors
$G_t=\{\mathbf g_{1,t},\ldots,\mathbf g_{M,t}\}$ with equal weights
$\mathbf v_t=(1/M)\mathbf 1$. Let
$\widetilde{\boldsymbol\Sigma}_{GG,t}$ be the posterior covariance matrix at
the anchors and $\widetilde{\boldsymbol\Sigma}_{\theta G,t}$ the posterior
cross-covariance between a candidate $\theta$ and the anchors. The average of
the GP function values at these anchors has posterior variance
\begin{equation}
 V_{G,t}=\mathbf v_t^{\top}
 \widetilde{\boldsymbol\Sigma}_{GG,t}\mathbf v_t
 \label{eq:region_variance}
\end{equation}
The covariance between the candidate value and this regional average is
$\widetilde{\boldsymbol\Sigma}_{\theta G,t}\mathbf v_t$. It determines the
candidate's adjustment:
\begin{equation}
 r_t(\theta)=
 \lambda_t\,
 \frac{\widetilde{\boldsymbol\Sigma}_{\theta G,t}\mathbf v_t}
 {\sqrt{V_{G,t}}},
 \quad
 \lambda_t\propto
 \widehat c_t\max\!\left(1-\frac{t-1}{T_{\mathrm L}},0\right)
 \label{eq:region_adjustment}
\end{equation}
Here, $\widehat c_t$ is the guidance score returned by the LLM, and
$T_{\mathrm L}$ is the guidance horizon. The coefficient $\lambda_t$ combines
this score with a fixed scale and a linearly decaying time factor. Positive posterior
cross-covariance produces a favorable adjustment for the minimization
problem; negative cross-covariance produces the opposite adjustment.

The adjustment $r_t(\theta)$ is rescaled using its magnitude at the regional
anchors and clipped to obtain a bounded value $\Delta_t(\theta)$. If
$s_{g,t}$ is the scale of the scalar GP target, EI is computed using
\begin{equation}
 \mu_t^{\mathrm{acq}}(\theta)
 =\mu_t(\theta)-s_{g,t}\Delta_t(\theta),
 \quad
 \sigma_t^{\mathrm{acq}}(\theta)=\sigma_t(\theta)
 \label{eq:region_acq_mean}
\end{equation}
Thus, a positive adjustment lowers the mean used by EI while leaving the
fitted GP and its posterior variance unchanged. Standard EI is used when no
valid region is available and after the guidance window.
Algorithm~\ref{alg:llmbo} summarizes the method.

\begin{algorithm}[!t]
\caption{LLMBO-MO}
\label{alg:llmbo}
\footnotesize
\begin{algorithmic}[1]
\Require Battery simulator, LLM, total budget $N$, initial budget $n_0$, weight set $W$, guidance horizon $T_{\mathrm L}$
\Ensure Feasible Pareto archive $P$
\State Generate and screen LLM-informed initial protocols
\State Combine selected LLM proposals with random valid protocols
\State Evaluate the $n_0$ initial protocols; initialize $D$ and $P$
\For{$t=1$ to $N-n_0$}
    \State Select $\mathbf w^{(t)}\in W$ and compute $g_t$ using Eq.~\eqref{eq:tchebycheff}
    \State Fit the scalar GP and construct EI
    \If{$t\le T_{\mathrm L}$}
        \State Query the LLM for a preferred region under $\mathbf w^{(t)}$
        \If{a valid region is returned}
            \State Apply Region-Lift using Eqs.~\eqref{eq:region_variance}--\eqref{eq:region_acq_mean}
        \EndIf
    \EndIf
    \State Maximize EI and select the next protocol $\theta_{t+1}$
    \State Evaluate $\theta_{t+1}$ and append the result to $D$
    \State Update the feasible Pareto archive $P$
\EndFor
\State \Return $P$
\end{algorithmic}
\end{algorithm}
